% --- Άσκηση 38861 (38861.tex) ---
\Askhsh{38861}{4}
\begin{ylh}{2}
    \item 2.1. Οι Πράξεις και οι Ιδιότητές τους
    \item 2.2. Διάταξη Πραγματικών Αριθμών
    \item 3.1. Εξισώσεις 1ου Βαθμού
    \item 3.3. Εξισώσεις 2ου Βαθμού
\end{ylh}
Υπάρχουν αρκετοί τύποι για τον υπολογισμό της δόσης ενός φαρμάκου για παιδιά με βάση τη δόση για ενήλικες. Κάποιοι από αυτούς χρησιμοποιούν την ηλικία του παιδιού και κάποιοι το σωματικό του βάρος. Οι τύποι που δίνονται στη συνέχεια χρησιμοποιούν την ηλικία του παιδιού και ισχύουν για παιδιά με μέσο βάρος και μέση ανάπτυξη.

Τύπος του Γιουνγκ (Young's Formula): Δόση για παιδιά$ = \dfrac{\varepsilon}{\varepsilon + 12} \cdot $Δόση για ενήλικες, όπου $\varepsilon$ είναι η ηλικία του παιδιού (σε έτη) από ενός έτους μέχρι 12 ετών.

Τύπος του Φράιντ (Fried's Formula): Δόση για παιδιά $ = \dfrac{\mu}{150} \cdot $Δόση για ενήλικες, όπου $\mu$ είναι η ηλικία του παιδιού (σε μήνες) έως 24 μηνών.
\begin{alist}
\item Αν το παιδί είναι $1,5$ έτους, να βρείτε τι μέρος της δόσης που αντιστοιχεί στον ενήλικα θα πάρει το παιδί, χρησιμοποιώντας καθέναν από τους δύο παραπάνω τύπους.
\monades{4}
\item
\begin{rlist}
\item Να βρείτε, χρησιμοποιώντας τον τύπο του Γιουνγκ, πόσων ετών πρέπει να είναι το παιδί, για να πάρει δόση φαρμάκου $60\si{mgr}$, εάν η δόση για τον ενήλικα είναι $300\si{mg}$.
\monades{5}
\item Να δείξετε ότι σύμφωνα με τον τύπο του Φράιντ, για τη δόση του παιδιού ισχύει:
$0 < \text{Δόση για παιδιά} \leq \dfrac{4}{25} \cdot \text{Δόση για ενήλικες}$.
\monades{5}
\item Θα μπορούσε με βάση τον τύπο του Φράιντ η δόση για το παιδί να είναι $60\si{mgr}$, αν η δόση για τον ενήλικα είναι $300\si{mgr}$; Να αιτιολογήσετε την απάντησή σας.
\monades{5}
\end{rlist}
\item Θα μπορούσε σε κάποια ηλικία το παιδί να πάρει την ίδια δόση του φαρμάκου με βάση και τους δύο παραπάνω τύπους; Να εξηγήσετε τη σκέψη σας.
\monades{6}

{\footnotesize Το παραπάνω θέμα αναπτύχθηκε στο πλαίσιο του έργου: «Ανάπτυξη Δοκιμασιών Αξιολόγησης Δεξιοτήτων Εγγραμματισμού στα μαθήματα της Νεοελληνικής Γλώσσας και Λογοτεχνίας, της Άλγεβρας, της Φυσικής και της Χημείας Α' Λυκείου Γενικού Λυκείου» Ανάδοχος: «Ειδικός Λογαριασμός Κονδυλίων Έρευνας (Ε.Λ.Κ.Ε) Πανεπιστημίου Ιωαννίνων» ΑΔΑΜ: 25SYMV016348911 2025-02-20.}
\end{alist}
